\makeatletter
\def\input@path{{./}{../}{../../}{../../../}{../../../../}{preamble/}{../preamble/}{../../preamble/}{../../../preamble/}{../../../../preamble/}}
\makeatother
\input{preamble_exams.tex}
\graphicspath{{./}{../}{../../}{../../../}{../../../../}{preamble/}{../preamble/}{../../preamble/}{../../../preamble/}{../../../../preamble/}}
\setlength{\columnsep}{6mm}
\linespread{1.0}
\setlength{\parskip}{3pt}
% Keep the shared heading's background inside the available column width.
\setlength{\fboxsep}{0pt}

\begin{document}
\ExamHeader{11th}{3}{Statistical models and research planning}{1st}\par
\vspace{2mm}
\noindent\textbf{Name:} \hrulefill\quad\textbf{Date:} \rule{3cm}{0.4pt}


Answer on separate paper and show your reasoning. Use the supplied normal and Student's t reference sheets.
Report Student's t critical values and normal critical values from the infinity row to \textbf{three decimal places}, and normal probabilities to \textbf{four decimal places}.
For Student's t cutoffs between printed entries, give probability bounds without interpolation.
Each criterion is worth \textbf{6 points} (\textbf{18 points} total); at least \textbf{4 points per criterion} are required to pass that criterion.

Questions 1--6 concern independent random observations from normal populations.
All critical values and probability questions refer to the sample mean standardized relative to the population mean; critical values have no units.

\begin{multicols}{2}
\raggedcolumns
\SubsectionBox{1. Export inspection costs \score{C8-\sml{2}}}
An exporter studies average inspection costs using twelve shipments. The population standard deviation is unknown and estimated from the sample.
\begin{enumerate}[label=\textbf{(\alph*)},leftmargin=*]
\item Determine the degrees of freedom and the positive critical value that leaves one percent of the probability to its right. \score{1}
\item Determine the negative critical value below which five percent of the probability lies. Justify its sign and the table column used. \score{1}
\end{enumerate}

\SubsectionBox{2. Market-stall electricity use \score{C8-\sml{2}}}
A study investigates average electricity use across twenty-five market stalls. The population standard deviation is unknown and estimated from the sample.
\begin{enumerate}[label=\textbf{(\alph*)},leftmargin=*]
\item Determine the degrees of freedom. For a central region containing ninety-eight percent of the probability, determine the combined probability outside it and the probability in each tail. \score{1}
\item Determine the two critical values, symmetric about zero, enclosing this central region. Identify the table column used. \score{1}
\end{enumerate}

\SubsectionBox{3. Cold-storage energy costs \score{C8-\sml{2}}}
An analyst studies average energy costs using fifteen cold-storage facilities. The population standard deviation is unknown and estimated from the sample.
\begin{enumerate}[label=\textbf{(\alph*)},leftmargin=*]
\item Determine the degrees of freedom and give bounds for the probability that the standardized sample mean exceeds positive 2.40. \score{1}
\item Give bounds for the probability that the standardized sample mean is below negative 2.40, and for the probability that it lies between negative 2.40 and positive 2.40. \score{1}
\end{enumerate}

\SubsectionBox{4. Bottling-line measurements \score{C9-\sml{2}}}
A bottling plant studies average bottle volume using sixteen bottles. The population standard deviation is known to be three millilitres.
\begin{enumerate}[label=\textbf{(\alph*)},leftmargin=*]
\item Choose and justify the distribution model for the standardized sample mean. Determine the negative critical value below which 2.5 percent of the probability lies. \score{1}
\item Calculate the probability that the standardized sample mean lies between negative 1.00 and positive 1.50. \score{1}
\end{enumerate}

\SubsectionBox{5. Repair-time averages \score{C9-\sml{2}}}
A workshop studies average repair time using ten repairs. The population standard deviation is unknown; the sample standard deviation is four minutes.
\begin{enumerate}[label=\textbf{(\alph*)},leftmargin=*]
\item Choose and justify the distribution model, determine the degrees of freedom, and find the positive critical value that the standardized sample mean exceeds with probability 2.5 percent. \score{1}
\item Estimate the probability that the standardized sample mean lies between negative 1.833 and positive 2.821. \score{1}
\end{enumerate}

\SubsectionBox{6. Comparing parcel-mass studies \score{C9-\sml{2}}}
Two studies investigate average parcel mass. Study A uses twelve parcels and a known population standard deviation of 0.4 kilograms. Study B uses twelve parcels; its population standard deviation is unknown and estimated by a sample standard deviation of 0.4 kilograms.
\begin{enumerate}[label=\textbf{(\alph*)},leftmargin=*]
\item Justify the distribution model for each study's standardized sample mean. Determine each model's two critical values, symmetric about zero, enclosing the central ninety-five percent of the probability. Calculate the difference between the positive critical values and explain it in terms of tail behaviour. \score{1}
\item Study B increases its sample to sixty-one parcels, with population standard deviation still unknown. State its model and degrees of freedom, and determine its new critical values enclosing the central ninety-five percent of the probability. Explain the trend toward the normal model and why sample size alone does not determine the model choice. \score{1}
\end{enumerate}

\SubsectionBox{7. Plan your future research \score{C10-\sml{6}}}
Plan three different real-life situations using actual observations:
a continuous variable modelled by a linear probability density over a stated range;
waiting times or intervals between events modelled by an exponential distribution;
and normally distributed measurements with an investigation of their population mean.
Address all three situations in each part. Submit a research plan only; no collected data, completed calculations, working links, or finished article are required.

\begin{enumerate}[label=\textbf{(\alph*)},leftmargin=*]
\item For each situation, specify the context, variable, units, population, time period, and research objective. For the normal situation, state a specific population-mean question. \score{2}

\item For each situation, identify a real public dataset and source, or give a feasible search strategy specifying a public source, search terms, variables, and observation coverage.
Describe planned data selection, organization, missing-value checks, and filters.
Explain how you would construct a histogram, choose its intervals, and calculate the sample mean, variance, and standard deviation.
How would you compare these with theoretical model statistics, obtain the required parameters, and address missing information?
For normal measurements, distinguish the distribution comparison from mean inference. \score{2}

\item Justify each proposed model and state the assumptions and evidence you will check.
For the normal situation, justify the planned normal or Student's t model for mean inference, addressing sampling, independence, and whether the population standard deviation is known or estimated; specify how degrees of freedom will be determined if applicable.
Identify one anticipated limitation or model mismatch for each situation.
Why are histograms alone insufficient to establish model suitability, and why must your choices remain provisional? \score{2}
\end{enumerate}
\end{multicols}
\end{document}

